% DOCUMENT INFO %%%%%%%%%%%%%%%%%%%%%%%%%%%%%%%%%%
\newcommand{\myDocType}{Report}  % eg, Research report, Preprint, blank...
\newcommand{\myAuthor}{%
  \styleAuthor Marco Coraggio\\
  \styleAffiliation v2.2 | Part of \href{https://www.marco-coraggio.com/scientist-s-matchbox}{Scientist's Matchbox}
  \and
  \styleAuthor Author 2\\
  \styleAffiliation Affiliation author 2
  \and
  \styleAuthor Author 3\\
  \styleAffiliation Affiliation author 3
  }
\newcommand{\myTitle}{Template for reports in \LaTeX}
\newcommand{\myHeaderAuthorTitle}{N.~Surname et al.: Title or short title if necessary} % leave blank to show nothing. If deleted, default is used: \myTitle
\newcommand{\fixedDate}{31 January 2099}  % used with appropriate class option


% PREAMBLE %%%%%%%%%%%%%%%%%%%%%%%%%%%%%%%%%%%%%%%
% Class ------------------------------------------
\documentclass[authorrefs]{lib/SMarticle}

% Options guide
% - all options in article class (new default: a4paper)
% - headers: looseheaders (default), thightheaders
% - date: todaydate (default), nodate, fixeddate (requires \fixedDate)
% - bibliography style: numberrefs (IEEEtran; default), authorrefs (IEEEtranSN). Note when changing from numberrefs to authorrefs, it may be necessary to run latexmk twice or delete the .bbl file.


% Front matter -----------------------------------
\begin{document}
\thanks{%
  This work was in part supported by\dots\\
  %
  N.~Surname is with the University of\dots (email@domain.it).
  N.~Surname is with the University of\dots (email@domain.it).\\
  \\
  \textcopyright\ 2019 Marco Coraggio, \href{https://creativecommons.org/licenses/by-nc-sa/4.0/}{CC BY-NC-SA} (Creative Commons Attribution-NonCommercial-ShareAlike 4.0 International License).
}
\maketitle
\thispagestyle{firstpage}  % apply page style for first page
\begin{abstract}
  In this document, \dots
\end{abstract}

\tableofcontents
%\clearpage  % uncomment to break page after table of contents


% BODY %%%%%%%%%%%%%%%%%%%%%%%%%%%%%%%%%%%%%%%%%%%
%-------------------------------------------------
\section{Section 1}%
\label{sec:section_1}

\cite{smith2024interesting}
\begin{theorem}[My theorem]
  \begin{equation}
    a = b
  \end{equation}
\end{theorem}
%
\begin{proof}
  Trivial!
\end{proof}

Imaginary unit with ``\texttt{\textbackslash imaginary}'': $\imaginary$ \par
Euler number with ``\texttt{\textbackslash euler}'': $\euler$ \par
Transpose with ``\texttt{\textbackslash T}'': $x\T$ \par
Upright math with ``\texttt{\textbackslash R\{\}}'': $\R{x}$ \par
Operator with ``\texttt{\textbackslash ON\{\}}'': $\ON{x}$ \par
Bold math with ``\texttt{\textbackslash B\{\}}'': $\B{x}$ \par
Mathbb with ``\texttt{\textbackslash BB\{\}}'': $\BB{X}$ \par
Mathcal with ``\texttt{\textbackslash C\{\}}'': $\C{X}$ \par
Absolute value with ``\texttt{\textbackslash abs\{\}}'': $\abs{x}$ \par
Norm with ``\texttt{\textbackslash norm\{\}}'': $\norm{x}$ \par


%-------------------------------------------------
\subsection{Subsection 1.1}%
\label{subsec:subsection_1_1}

Important equation: \boxMath{x = a + b}.

\paragraph{Paragraph} Text.

\subparagraph{Subparagraph} Text.


\begin{figure}[t]
   \centering
   \href{https://www.marco-coraggio.com/scientist-s-matchbox}{\includegraphics[max width=0.7\columnwidth]{logo_title_scientist's_matchbox.png}}
   \caption{My caption.}
   \label{fig:fig_label}
\end{figure}

\begin{figure}[t]
   \centering
   \subfloat[subcaption]{\includegraphics[max width=0.49\columnwidth]{example-image-a}
       \label{fig:fig_label_2}}
   \subfloat[subcaption]{\includegraphics[max width=0.49\columnwidth]{example-image-b}
       \label{fig:fig_label_3}}
   \caption{(a) My caption. (b) My caption.}
   \label{fig:whole_figure}
\end{figure}

\begin{table}[t]
  \centering
  \begin{tabular}{l|lll}
      \toprule
      Column 1  &  Column 2  &  Column 3  &  Column 4  \\
      \midrule
      \multicolumn{4}{c}{\textit{Group 1}}\\
      \noalign{\smallskip}
      Row 1  &  value  &  value  &  value  \\
      Row 2  &  value  &  value  &  value  \\
      \midrule
      \multicolumn{4}{c}{\textit{Group 2}}\\
      \noalign{\smallskip}
      Row 1  &  value  &  value  &  value  \\
      Row 2  &  value  &  value  &  value  \\
      \bottomrule
  \end{tabular}
  \caption{My caption.}
  \label{tab:tab_label}
  \end{table}

\begin{algorithm}[t]
  \caption{Name of algorithm}
  \label{alg:algorithm_label}
  %
  %\DontPrintSemicolon  % to avoid printing semicolons
  \SetKwProg{Fn}{Function}{}{}
  \SetKwInput{KwParam}{Parameters}
  \SetKwFunction{myKwFun}{fun\textunderscore name}
  %
  \KwIn{input 1; input 2.}%  
  \KwOut{output 1; output 2.}
  \KwParam{param 1; param 2.}
  %
  \BlankLine
  %
  $a \leftarrow \texttt{fun\textunderscore name}(1)$\; \label{line:a}% close commands with \; 
  $b \leftarrow 2$%
    \tcp*{a comment} % this comment automatically closes line
  \While(\tcp*[f]{a comment for a line without end}){$a > 0$}{
    $c \leftarrow 3$\;
    \lIf{$a <1 $}{something, referencing instruction \ref{line:a}} % l makes loop in one line
  }
  \BlankLine
  \setcounter{AlgoLine}{0}
  %
  \Fn{\myKwFun{myargs}}{
      do something\;
  }
  % more info: http://mirror.ox.ac.uk/sites/ctan.org/macros/latex/contrib/algorithm2e/doc/algorithm2e.pdf
\end{algorithm}


% REFERENCES %%%%%%%%%%%%%%%%%%%%%%%%%%%%%%%%%%%%%
\bibliographystyle{\bibliographyStyleName}
\bibliography{res/references.bib}

\end{document}



%%%%%%%%%%%%%%%%%%%%%%%%%%%%%%%%%%%%%%%%%%%%%%%%%%
% RESOURCES %%%%%%%%%%%%%%%%%%%%%%%%%%%%%%%%%%%%%%
%%%%%%%%%%%%%%%%%%%%%%%%%%%%%%%%%%%%%%%%%%%%%%%%%%


% Appendices and acknowledgment ------------------
% do not use \section anymore after \appendix, only \section* is possibly needed. If single appendix:
%\appendix[Proof of the Zonklar Equations]
% or
%\appendix  % for no appendix heading

% If more than one appendix
%\appendices
%\section{Proof of the First Zonklar Equation}
%\section{}

% use section* for acknowledgment
%\section*{Acknowledgment}
%The authors would like to thank...
%-------------------------------------------------